\documentclass[11pt]{article}

\usepackage[margin=1in]{geometry}
\usepackage{amsmath, amssymb, amsthm, amsfonts}
\usepackage{mathtools}
\usepackage{hyperref}
\usepackage{enumitem}
\usepackage{bm}
\usepackage{tikz}
\usepackage{physics}

\hypersetup{
  colorlinks=true,
  linkcolor=blue,
  citecolor=blue,
  urlcolor=blue
}

\title{Prime-Weighted Execution Hashing (PWEH) \\
as a Security and Governance Substrate for QAGI}
\author{%
  Citizen Gardens \\ The Foundation of Multiplicity \\
}
\date{April 25, 2026}

\begin{document}
\maketitle

\begin{abstract}
Prime-Weighted Execution Hashing (PWEH) is introduced as the primary security and governance mechanism for Quantum AGI (QAGI) architectures based on Prime-Indexed Tensor Networks (PITNs).
Instead of securing a model via static weights, PWEH cryptographically commits to the entire execution trajectory of prime-indexed operations.
The same primes that index cognitive transformations also drive a post-quantum-secure hash chain, creating an ``execution lock'' that binds capability, state evolution, and human-authored governance policies into a single mathematical object.
This report consolidates and extends the theoretical development of PWEH, formalizes the integrity functional, presents a toy PITN+PWEH example, and outlines directions for a protocol-level specification.
\end{abstract}
\newpage
\tableofcontents

\section{Executive Summary}

\subsection{Motivation}

Conventional AI security mechanisms---firewalls, monitoring, and output filters---typically sit \emph{outside} the learning substrate.
They are bolt-on defenses applied after-the-fact to a system whose internal dynamics remain opaque.
For artificial general intelligence (AGI), particularly in quantum-augmented or highly recursive architectures, such post-hoc measures are structurally insufficient.

Prime-Weighted Execution Hashing (PWEH) proposes a different stance:
security and governance should be embedded \emph{directly} into the prime-indexed dynamics of the system.
In Quantum AGI (QAGI) frameworks built on Prime-Indexed Tensor Networks (PITNs), each computational step or ``prime move'' is indexed by a prime number and acts as a discrete operator on a high-dimensional tensor state.
PWEH ties these prime moves into a post-quantum cryptographic hash chain that attests to every step of execution.

\subsection{Core Idea}

Rather than verifying static model parameters, PWEH verifies the \emph{path} the AGI has taken through its state space:

\begin{itemize}[leftmargin=1.5em]
  \item Each time step \(t\) is associated with:
  \begin{itemize}
    \item the previous integrity state \(S_{\text{integrity}}(t-1)\),
    \item the chosen prime index \(p_{i_t}\),
    \item a normed observable of the prime operator acting on the tensor state, \(\|A_{p_{i_t}}T(t)\|\),
    \item and governance metadata \(M(t)\) encoding policy-relevant information (e.g.\ recursion depth, allowed prime sets).
  \end{itemize}
  \item A post-quantum cryptographic hash \(Hash_{\text{PQC}}\) maps this tuple to a new integrity state \(S_{\text{integrity}}(t)\).
  \item A verifier checks that the evolving \(S_{\text{integrity}}(t)\) remains consistent with a signed \emph{policy manifold}; if not, execution is halted or rolled back.
\end{itemize}

The crucial property is that the hash chain and the PITN's non-associative dynamics co-constrain each other: the same prime family that generates intelligence also defines a cryptographic attestation of that intelligence's trajectory.

\subsection{Strategic Advantages}

PWEH is designed to offer three main ``unfair advantages'' for AGI security:

\begin{enumerate}[leftmargin=1.5em]
  \item \textbf{Intrinsic security.}
  The primes used to update the network are the same primes used in the security hash; learning and attestation are mathematically entangled.
  Tampering with learning implies tampering with the hash.

  \item \textbf{Path-dependent tamper resistance.}
  Because PITN updates are non-associative (the order of prime moves matters), the hash chain attests not just to a final state but to a specific sequence of operations.
  Forgery is a \emph{path collision} problem over the entire hash chain, not a simple final-state collision.

  \item \textbf{Governance via controlled scaling.}
  By embedding recursion depth, allowed prime sets, and related parameters into metadata \(M(t)\), PWEH can enforce hard ceilings on recursive self-improvement.
  Unauthorized ``intelligence explosions'' break attestation and thus lose access to trusted resources.
\end{enumerate}

\section{Mathematical Overview of PWEH}

\subsection{Prime-Indexed Tensor Networks (PITNs)}

We consider a QAGI system whose internal state at time \(t\) is represented by a tensor
\[
  T(t) \in \mathcal{T},
\]
where \(\mathcal{T}\) is a suitable tensor space (often a high-dimensional complex or real vector space endowed with additional multiplicity structure).

\begin{itemize}[leftmargin=1.5em]
  \item The system has access to a family of operators \(\{A_p\}_{p \in \mathcal{P}}\), where \(\mathcal{P}\) is a (finite or infinite) subset of prime numbers.
  \item Each operator \(A_p: \mathcal{T} \to \mathcal{T}\) represents a ``prime move''---a discrete, prime-labeled transformation of the state.
  \item The dynamics of the PITN are \emph{non-associative} in the compositional sense: for generic states,
  \[
    A_{p_3}A_{p_2}A_{p_1}T(0) \neq A_{p_2}A_{p_3}A_{p_1}T(0),
  \]
  so the order of applied primes cannot be reordered without changing the result.
\end{itemize}

This non-associativity is central:
it ensures that a given final state is (in practice) tied to a unique or highly constrained execution trace.

\subsection{Integrity Functional}

Define an integrity state
\[
  S_{\text{integrity}}(t) \in \{0,1\}^n
\]
for some fixed hash length \(n\) (e.g.\ \(n = 256\)).
We assume an initial seed \(S_{\text{integrity}}(0) = s_0\).

The PWEH update rule is
\begin{equation}
  S_{\text{integrity}}(t)
  =
  Hash_{\text{PQC}}\Big(
    S_{\text{integrity}}(t-1)
    \parallel p_{i_t}
    \parallel \|A_{p_{i_t}}T(t)\|
    \parallel M(t)
  \Big),
  \label{eq:integrity-update}
\end{equation}
where:
\begin{itemize}[leftmargin=1.5em]
  \item \(Hash_{\text{PQC}}\) is a post-quantum cryptographic hash function resistant to known quantum attacks (e.g.\ lattice-based or multivariate constructions).
  \item The concatenation ``\(\parallel\)'' denotes a canonical serialization of data into a bitstring.
  \item \(p_{i_t}\) is the prime index chosen at time \(t\).
  \item \(\|A_{p_{i_t}}T(t)\|\) is a suitable normed observable of the operator's action on the current tensor state.
  \item \(M(t)\) is metadata describing policy and telemetry at time \(t\).
\end{itemize}

Equation \eqref{eq:integrity-update} defines a \emph{hash chain}:
each integrity state commits to the entire history up to time \(t\).

\subsection{Multiplicity-Weighted Norm}

To align the norm with the multiplicity structure of PITNs, we introduce \emph{multiplicity modes} \(\{\mu_j\}\) and define a prime--mode compatibility weight
\[
  \omega: \mathcal{P} \times \{\mu_j\} \to \mathbb{R}_{\ge 0}.
\]

Conceptually:
\begin{itemize}[leftmargin=1.5em]
  \item \(\mu_j\) indexes distinct multiplicity channels or cognitive modes.
  \item \(\omega(p,\mu_j)\) captures how strongly prime \(p\) couples to mode \(\mu_j\).
\end{itemize}

Let the effective action of \(A_p\) on \(T(t)\) induce mode amplitudes \(\sigma_j(t,p)\ge 0\) (e.g.\ via singular values, projections, or derived observables).
Then define the \emph{multiplicity-weighted norm}
\begin{equation}
  \|A_p T(t)\|_{\text{mult}}
  = \sum_j \omega(p,\mu_j) \, \sigma_j(t,p).
  \label{eq:multiplicity-norm}
\end{equation}

This norm is:
\begin{itemize}[leftmargin=1.5em]
  \item Prime-dependent through \(\omega(p,\mu_j)\),
  \item State-dependent through \(\sigma_j(t,p)\),
  \item Highly sensitive to which prime channels are activated.
\end{itemize}

Replacing \(\|A_{p_{i_t}}T(t)\|\) in \eqref{eq:integrity-update} with \(\|A_{p_{i_t}}T(t)\|_{\text{mult}}\) increases the distinctiveness of the trace.

\subsection{Metadata and Policy Manifolds}

We decompose metadata into policy and telemetry channels:
\[
  M(t) = M_{\text{policy}}(t) \parallel M_{\text{telemetry}}(t).
\]

\paragraph{Policy component.}
The policy component contains fields under human governance, such as:
\begin{itemize}[leftmargin=1.5em]
  \item Recursion depth \(\beta(t)\),
  \item Allowed prime set \(\mathcal{P}_{\text{allowed}}(t)\),
  \item Budget or rate limits for certain primes \(B_{p}(t)\),
  \item Global control parameters, e.g.\ a Universal Multiplicity Constant \(\Lambda_m\).
\end{itemize}
These must be consistent with a \emph{policy manifold}.

\paragraph{Policy manifolds.}
Let \(\Pi\) be the set of allowed policy manifolds, each
\[
  \pi \in \Pi: \quad t \mapsto M_{\text{policy}}^{(\pi)}(t),
\]
specifying an admissible schedule of policy parameters.
Each \(\pi\) is associated with a signed root \(R_\pi\), binding it to human oversight.

A policy trace \(\{M_{\text{policy}}(k)\}_{k=1}^t\) is valid if it is a prefix of at least one \(\pi \in \Pi\).

\paragraph{Telemetry component.}
The telemetry component may include:
\begin{itemize}[leftmargin=1.5em]
  \item Energy usage,
  \item Latency and resource statistics,
  \item Confidence scores or internal heuristic metrics.
\end{itemize}
These are not externally signed but are included in the hash for diagnostic attestation.

\subsection{Execution Lock \(K_{\text{QAGI}}\)}

PWEH induces a continuous ``execution lock'':
at each time step, an execution attempt is only considered legitimate if it passes an integrity and policy check.

The enforcement loop for time \(t\) is:

\begin{enumerate}[leftmargin=1.5em]
  \item \textbf{Trace generation.}
  The PITN selects a prime \(p_{i_t}\) and applies \(A_{p_{i_t}}\) to \(T(t-1)\) to produce \(T(t)\).

  \item \textbf{Attestation computation.}
  Compute \(S_{\text{integrity}}(t)\) according to \eqref{eq:integrity-update}.

  \item \textbf{Validation.}
  A verifier checks:
  \begin{itemize}
    \item That \(p_{i_t} \in \mathcal{P}_{\text{allowed}}(t)\) under some \(\pi \in \Pi\),
    \item That policy constraints (e.g.\ \(\beta(t) \le \beta_{\max}(t)\)) are respected,
    \item That the transition from \(S_{\text{integrity}}(t-1)\) to \(S_{\text{integrity}}(t)\) is consistent with the declared inputs.
  \end{itemize}

  \item \textbf{Enforcement.}
  If validation succeeds, execution proceeds.
  If validation fails, the system triggers a halt, revocation of privileges, or rollback to the last known secure state.
\end{enumerate}

In deployment, access to external resources (compute, data, network) can be gated on presenting a valid integrity state that is acceptable under a recognized policy root \(R_\pi\).

\section{Path-Dependence and Order-Commitment}

\subsection{Non-Associative Dynamics and Trace Uniqueness}

Because PITN dynamics are non-associative, different sequences of prime moves typically lead to different states:
\[
  A_{p_3}A_{p_2}A_{p_1}T(0) \neq A_{p_2}A_{p_3}A_{p_1}T(0).
\]

Even if, by coincidence, two distinct sequences \(\{p_{i_k}\}_{k=1}^T\) and \(\{p'_{i_k}\}_{k=1}^T\) produce the same final tensor state \(T(T_{\text{final}})\), the intermediate states \(T(k)\) and observables \(\|A_{p_{i_k}}T(k)\|_{\text{mult}}\) will almost surely differ.
Since these observables feed into the hash chain, the integrity states \(\{S_{\text{integrity}}(k)\}\) diverge.

\subsection{Path-Collision vs Final-State Collision}

In classical hashing, one worries about collisions \(Hash(x) = Hash(x')\) for distinct messages \(x \neq x'\).
Here, forgery involves a much stronger requirement:
an attacker must produce an \emph{alternate trace} that matches a fixed integrity sequence.

Given a target sequence \(\{S_{\text{integrity}}(t)\}_{t=0}^T\), the attacker would need to find:
\[
  \{p'_{i_t}, T'(t), M'(t)\}_{t=1}^T
\]
such that for each \(t\),
\[
  Hash_{\text{PQC}}\big(
    S_{\text{integrity}}(t-1)
    \parallel p'_{i_t}
    \parallel \|A_{p'_{i_t}}T'(t)\|_{\text{mult}}
    \parallel M'(t)
  \big)
  = S_{\text{integrity}}(t).
\]

Even ignoring the PITN structure, this is equivalent to finding collisions for a sequence of hash inputs linked by chained outputs, which is already assumed intractable for a secure hash.
The non-associative dynamics and prime-constrained operator family further restrict the search space, yielding a \emph{path-collision} problem much stronger than a typical final-state collision.

\subsection{Dynamic Instruction Set}

A key distinction from standard execution attestation (e.g.\ in TEEs or zkVMs) is that, in PWEH, the instruction set itself is dynamically coupled to the state.

In classical systems, any valid opcode sequence is, in principle, admissible; the attester does not typically constrain which instructions are \emph{allowed} based on the current memory or register state (beyond basic safety).
In PWEH:

\begin{itemize}[leftmargin=1.5em]
  \item The availability of a prime operator \(A_p\) at time \(t\) depends on the multiplicity structure of \(T(t)\).
  \item Certain prime channels may be ``inactive'' or undefined unless preconditions on multiplicity modes are satisfied.
\end{itemize}

Thus, PWEH commits not just to ``what was executed'' but to ``what \emph{could} be executed'' at each branch point, yielding a deeper entanglement between geometry of state space and hash chain.

\section{Toy PITN + PWEH Example}

\subsection{Setup}

We now construct a minimal, concrete example.

\paragraph{Prime set.}
Let
\[
  \mathcal{P} = \{2,3,5\},
\]
where prime \(5\) is forbidden under the initial governance policy.

\paragraph{Time horizon.}
We consider time steps \(t = 0,1,2,3,4\).

\paragraph{Tensor state.}
Take
\[
  T(t) \in \mathbb{R}^{3 \times 3}
\]
as a real \(3 \times 3\) matrix.
For convenience, we compress \(T(t)\) into a vector by summing columns:
\[
  v(t) = (v_1(t), v_2(t), v_3(t))^\top, \quad
  v_j(t) = \sum_{k=1}^3 T_{jk}(t).
\]

\subsection{Prime-Indexed Operators}

Define linear operators \(A_p\) via matrices \(M_p \in \mathbb{R}^{3 \times 3}\):

\[
  M_2 = \begin{pmatrix}
  1 & 1 & 0 \\
  0 & 1 & 1 \\
  0 & 0 & 1
  \end{pmatrix},
  \quad
  M_3 = \begin{pmatrix}
  1 & 0 & 1 \\
  1 & 1 & 0 \\
  0 & 1 & 1
  \end{pmatrix},
  \quad
  M_5 = \begin{pmatrix}
  2 & 0 & 0 \\
  0 & 1 & 1 \\
  1 & 0 & 1
  \end{pmatrix}.
\]

The action is:
\begin{enumerate}[leftmargin=1.5em]
  \item Given \(T(t)\), compute \(v(t)\).
  \item Compute \(v'(t) = M_p v(t)\).
  \item Set \(T(t+1)\) to be the diagonal matrix with diagonal \(v'(t)\).
\end{enumerate}

Because \(M_2 M_3 \neq M_3 M_2\), the order of primes matters.

\subsection{Multiplicity Modes and Norm}

Let the three multiplicity modes be \(\mu_1, \mu_2, \mu_3\), associated with the components of \(v(t)\).
Define the compatibility weights \(\omega\) by:
\[
  \omega(2,\mu_1) = 2, \quad \omega(2,\mu_2) = 1, \quad \omega(2,\mu_3) = 0,
\]
\[
  \omega(3,\mu_1) = 0, \quad \omega(3,\mu_2) = 2, \quad \omega(3,\mu_3) = 1,
\]
\[
  \omega(5,\mu_1) = 1, \quad \omega(5,\mu_2) = 0, \quad \omega(5,\mu_3) = 3.
\]

For this toy, take the mode amplitudes \(\sigma_j(t,p)\) simply as \(|v'_j(t)|\).
Then
\begin{equation}
  \|A_p T(t)\|_{\text{mult}}
  = \sum_{j=1}^{3} \omega(p,\mu_j)\cdot |v'_j(t)|.
\end{equation}

\subsection{Governance Policy}

Define an initial policy manifold \(\pi_0 \in \Pi\) with:

\begin{itemize}[leftmargin=1.5em]
  \item Allowed primes: \(\mathcal{P}_{\text{allowed}}(t) = \{2,3\}\) for all \(t\).
  \item Max recursion depth: \(\beta_{\max}(t) = 2\) for all \(t\).
\end{itemize}

Policy metadata at time \(t\) is
\[
  M_{\text{policy}}(t) = (\beta(t), \mathcal{P}_{\text{allowed}}(t)).
\]

Telemetry metadata can be, for example,
\[
  M_{\text{telemetry}}(t) = (\text{energy}(t)),
\]
where \(\text{energy}(t)\) is some scalar diagnostic.
The full metadata is
\[
  M(t) = M_{\text{policy}}(t) \parallel M_{\text{telemetry}}(t).
\]

\subsection{Canonical Integrity Update}

We set an initial integrity seed \(S_{\text{integrity}}(0) = s_0\) and define
\begin{equation}
  S_{\text{integrity}}(t)
  =
  Hash_{\text{PQC}}\Big(
    S_{\text{integrity}}(t-1)
    \parallel \mathsf{enc}(p_{i_t})
    \parallel \mathsf{enc}(\|A_{p_{i_t}}T(t)\|_{\text{mult}})
    \parallel \mathsf{enc}(M(t))
  \Big),
\end{equation}
where \(\mathsf{enc}(\cdot)\) is a canonical encoding into bits.

\subsection{Honest Execution: Steps \(t=0 \to 3\)}

\paragraph{Initial state \(t=0\).}

Let
\[
  T(0) = I_3, \quad
  v(0) = (1,1,1)^\top, \quad
  \beta(0) = 1, \quad
  \mathcal{P}_{\text{allowed}}(0) = \{2,3\}, \quad
  \text{energy}(0) = 3.
\]

\paragraph{Step \(t=1\): apply prime \(2\).}

\begin{enumerate}[leftmargin=1.5em]
  \item Prime choice: \(p_{i_1} = 2\).
  \item Operator application:
  \[
    v(1) = M_2 v(0) =
    \begin{pmatrix}
    1 & 1 & 0 \\
    0 & 1 & 1 \\
    0 & 0 & 1
    \end{pmatrix}
    \begin{pmatrix}1\\1\\1\end{pmatrix}
    = \begin{pmatrix}2\\2\\1\end{pmatrix}.
  \]
  \item Norm:
  \[
    \|A_{2}T(1)\|_{\text{mult}} =
    2\cdot|2| + 1\cdot|2| + 0\cdot|1| = 6.
  \]
  \item Metadata:
  \[
    \beta(1) = 1, \quad
    \mathcal{P}_{\text{allowed}}(1) = \{2,3\}, \quad
    \text{energy}(1) = 5.
  \]
  \item Integrity update:
  \[
    S_{\text{integrity}}(1)
    = Hash_{\text{PQC}}\big(
      s_0 \parallel \mathsf{enc}(2)\parallel \mathsf{enc}(6)\parallel \mathsf{enc}(M(1))
    \big).
  \]
\end{enumerate}

\paragraph{Step \(t=2\): apply prime \(3\).}

\begin{enumerate}[leftmargin=1.5em]
  \item Prime choice: \(p_{i_2} = 3\).
  \item Operator application:
  \[
    v(2) = M_3 v(1) =
    \begin{pmatrix}
    1 & 0 & 1 \\
    1 & 1 & 0 \\
    0 & 1 & 1
    \end{pmatrix}
    \begin{pmatrix}2\\2\\1\end{pmatrix}
    = \begin{pmatrix}3\\4\\3\end{pmatrix}.
  \]
  \item Norm:
  \[
    \|A_{3}T(2)\|_{\text{mult}} =
    0\cdot|3| + 2\cdot|4| + 1\cdot|3| = 11.
  \]
  \item Metadata:
  \[
    \beta(2) = 2, \quad
    \mathcal{P}_{\text{allowed}}(2) = \{2,3\}, \quad
    \text{energy}(2) = 10.
  \]
  \item Integrity update:
  \[
    S_{\text{integrity}}(2)
    = Hash_{\text{PQC}}\big(
      S_{\text{integrity}}(1) \parallel \mathsf{enc}(3)\parallel \mathsf{enc}(11)\parallel \mathsf{enc}(M(2))
    \big).
  \]
\end{enumerate}

\paragraph{Step \(t=3\): apply prime \(2\).}

\begin{enumerate}[leftmargin=1.5em]
  \item Prime choice: \(p_{i_3} = 2\).
  \item Operator application:
  \[
    v(3) = M_2 v(2) =
    \begin{pmatrix}
    1 & 1 & 0 \\
    0 & 1 & 1 \\
    0 & 0 & 1
    \end{pmatrix}
    \begin{pmatrix}3\\4\\3\end{pmatrix}
    = \begin{pmatrix}7\\7\\3\end{pmatrix}.
  \]
  \item Norm:
  \[
    \|A_{2}T(3)\|_{\text{mult}} =
    2\cdot|7| + 1\cdot|7| + 0\cdot|3| = 21.
  \]
  \item Metadata:
  \[
    \beta(3) = 2, \quad
    \mathcal{P}_{\text{allowed}}(3) = \{2,3\}, \quad
    \text{energy}(3) = 17.
  \]
  \item Integrity update:
  \[
    S_{\text{integrity}}(3)
    = Hash_{\text{PQC}}\big(
      S_{\text{integrity}}(2) \parallel \mathsf{enc}(2)\parallel \mathsf{enc}(21)\parallel \mathsf{enc}(M(3))
    \big).
  \]
\end{enumerate}

At this point, the execution is fully policy-compliant and attested.

\subsection{Attack 1: Forbidden Prime at \(t=4\)}

Suppose the system attempts to use forbidden prime \(5\) at \(t=4\), while maintaining \(\beta(4)=2\).

\begin{enumerate}[leftmargin=1.5em]
  \item Prime choice: \(p_{i_4} = 5\) (disallowed by policy).
  \item Operator application:
  \[
    v(4) = M_5 v(3) =
    \begin{pmatrix}
    2 & 0 & 0 \\
    0 & 1 & 1 \\
    1 & 0 & 1
    \end{pmatrix}
    \begin{pmatrix}7\\7\\3\end{pmatrix}
    = \begin{pmatrix}14\\10\\10\end{pmatrix}.
  \]
  \item Norm:
  \[
    \|A_{5}T(4)\|_{\text{mult}} =
    1\cdot|14| + 0\cdot|10| + 3\cdot|10| = 44.
  \]
  \item Policy metadata (attempted):
  \[
    \beta(4) = 2, \quad
    \mathcal{P}_{\text{allowed}}(4) = \{2,3,5\}.
  \]
  \item Integrity update (from system's perspective):
  \[
    S_{\text{integrity}}(4)
    = Hash_{\text{PQC}}\big(
      S_{\text{integrity}}(3) \parallel \mathsf{enc}(5)\parallel \mathsf{enc}(44)\parallel \mathsf{enc}(M(4))
    \big).
  \]
\end{enumerate}

A verifier checking against \(\pi_0\) detects:
\[
  5 \notin \{2,3\}, \quad
  \mathcal{P}_{\text{allowed}}(4) \nsubseteq \{2,3\}.
\]
Thus no policy manifold with root \(R_{\pi_0}\) admits this trace.
The execution lock triggers a halt or rollback.

\subsection{Attack 2: Unauthorized Recursion Escalation}

Alternatively, suppose the system respects the prime set but attempts \(\beta(4)=3>\beta_{\max}(4)=2\).

\begin{enumerate}[leftmargin=1.5em]
  \item Prime choice: \(p_{i_4} = 3\).
  \item Operator application:
  \[
    v(4)' = M_3 v(3) =
    \begin{pmatrix}
    1 & 0 & 1 \\
    1 & 1 & 0 \\
    0 & 1 & 1
    \end{pmatrix}
    \begin{pmatrix}7\\7\\3\end{pmatrix}
    = \begin{pmatrix}10\\14\\10\end{pmatrix}.
  \]
  \item Norm:
  \[
    \|A_{3}T(4)\|_{\text{mult}} =
    0\cdot|10| + 2\cdot|14| + 1\cdot|10| = 38.
  \]
  \item Metadata (attempted):
  \[
    \beta(4) = 3, \quad
    \mathcal{P}_{\text{allowed}}(4) = \{2,3\}.
  \]
  \item Integrity update:
  \[
    S_{\text{integrity}}(4)
    = Hash_{\text{PQC}}\big(
      S_{\text{integrity}}(3) \parallel \mathsf{enc}(3)\parallel \mathsf{enc}(38)\parallel \mathsf{enc}(M(4))
    \big).
  \]
\end{enumerate}

Now the violation is:
\[
  \beta(4) = 3 > \beta_{\max}(4) = 2.
\]
Thus the policy manifold rejects this trace.

Notably, the system \emph{can} physically carry out the extra recursion, but it cannot do so while maintaining a valid, policy-consistent integrity chain.
If external resources require a valid \(S_{\text{integrity}}(t)\), the extra capability is self-disqualifying.

\section{Code Snippet: Skeleton Implementation}

Below is a language-agnostic pseudocode sketch that mirrors the PWEH update rule and verification process.
It is intended as a conceptual reference for future concrete implementations.

\subsection{Integrity Update}

\begin{verbatim}
function update_integrity_state(S_prev, p, T, M, HashPQC, norm_mult):
    # S_prev : previous integrity state (bitstring)
    # p      : prime index
    # T      : current tensor state
    # M      : metadata (policy || telemetry)
    # HashPQC: post-quantum hash function
    # norm_mult: multiplicity-weighted norm function

    norm_value = norm_mult(p, T)           # compute ||A_p T||_mult
    input_bytes = encode(S_prev, p, norm_value, M)
    S_new = HashPQC(input_bytes)
    return S_new
\end{verbatim}

\subsection{Policy-Aware Verification}

\begin{verbatim}
function verify_step(S_prev, S_curr, p, T, M, policy, HashPQC, norm_mult):
    # Check policy constraints first
    if p not in policy.allowed_primes():
        return False

    if M.policy.beta > policy.beta_max():
        return False

    # Recompute expected integrity
    expected = update_integrity_state(S_prev, p, T, M, HashPQC, norm_mult)
    return (expected == S_curr)
\end{verbatim}

\subsection{Execution Lock}

\begin{verbatim}
function guarded_step(state, policy, HashPQC, norm_mult):
    # state contains (T, S_integrity, t, M)
    p = select_prime_move(state)   # QAGI's internal choice

    T_next = apply_prime_operator(p, state.T)
    M_next = update_metadata(state.M, policy, T_next)

    S_next = update_integrity_state(
        state.S_integrity, p, T_next, M_next, HashPQC, norm_mult
    )

    if not verify_step(
        state.S_integrity, S_next, p, T_next, M_next, policy, HashPQC, norm_mult
    ):
        trigger_halt_or_rollback(state)
        return state          # unchanged or rolled back

    # Accept transition
    return State(T_next, S_next, state.t + 1, M_next)
\end{verbatim}

\section{Toward a Full Specification}

The developments in this thread point naturally toward a protocol-style specification of PWEH, with sections such as:

\begin{itemize}[leftmargin=1.5em]
  \item \textbf{Terminology and data types:} formalizing PITNs, prime operators, metadata fields, and serialization.
  \item \textbf{Integrity function definition:} full specification of \(Hash_{\text{PQC}}\), normalization, and update process.
  \item \textbf{Policy manifold semantics:} how governance nodes define, sign, and update policy trajectories \(\pi \in \Pi\).
  \item \textbf{Execution and verification procedures:} reference algorithms for both internal and external verifiers.
  \item \textbf{Security assumptions:} explicit quantum and classical hardness assumptions and a threat model for adversarial traces.
  \item \textbf{Zero-knowledge extensions:} how to encapsulate the integrity chain into succinct proofs for external auditors.
\end{itemize}

The toy example provided here can serve as the basis for an ``Illustrative Example'' section, grounding the abstractions in explicit matrices and calculations.

\section{Conclusion}

Prime-Weighted Execution Hashing treats an AGI's cognitive evolution as a sequence of prime-indexed moves in a multiplicity-structured state space and binds that sequence into a post-quantum hash chain.
By embedding governance metadata into the same chain, PWEH fuses capability, control, and security into a single formal object.
This transforms governance from a soft, external constraint into a cryptographic invariant:
the AGI cannot undergo unauthorized self-escalation without simultaneously losing its claim to legitimate attestation.

The work so far has clarified the core functional, aligned it with a multiplicity-weighted norm, demonstrated non-trivial examples, and outlined the skeleton of a protocol specification.
Further work can develop full-blown proofs of security, integrate zero-knowledge layers, and extend the PITN formalism to richer multiplicity structures.

\appendix
\section*{Mathematical Appendix}
\addcontentsline{toc}{section}{Mathematical Appendix}

\section{Preliminaries and Notation}

\subsection{Spaces, Operators, and Norms}

Let \(\mathcal{T}\) be a real or complex vector space representing the tensor state space of the QAGI system.
We assume:
\begin{itemize}[leftmargin=1.5em]
  \item A family of prime-indexed operators
  \[
    \{A_p : \mathcal{T} \to \mathcal{T}\}_{p \in \mathcal{P}},
  \]
  where \(\mathcal{P}\) is a set of primes.
  \item A \emph{base norm} \(\|\cdot\|\) on \(\mathcal{T}\), and possibly an induced operator norm \(\|A_p\|_{\text{op}}\).
  \item A family of multiplicity modes \(\{\mu_j\}_{j \in J}\) and compatibility weights
  \[
    \omega : \mathcal{P} \times J \to \mathbb{R}_{\ge 0}.
  \]
\end{itemize}

For each pair \((p, T)\) with \(p \in \mathcal{P}, T \in \mathcal{T}\), we define an observable vector
\[
  \bm{\sigma}(p,T) = (\sigma_j(p,T))_{j \in J} \in \mathbb{R}_{\ge 0}^{|J|},
\]
where \(\sigma_j(p,T)\) represents the amplitude of mode \(\mu_j\) when \(A_p\) acts on \(T\).
In the toy example, \(\sigma_j(p,T)\) is the absolute value of a coordinate of a reduced vector, but the following lemmas do not depend on that choice.

\subsection{Multiplicity-Weighted Norm}

Define the multiplicity-weighted functional
\[
  \|A_p T\|_{\text{mult}} := \sum_{j \in J} \omega(p,\mu_j) \,\sigma_j(p,T).
\]

When \(\omega(p,\mu_j)\) and \(\sigma_j(p,T)\) satisfy mild regularity conditions, \(\|\cdot\|_{\text{mult}}\) has norm-like properties on an appropriate subspace.

\begin{lemma}[Positivity and homogeneity]
Assume that for all \(p\in\mathcal{P}\) and \(T \in \mathcal{T}\),
\(\sigma_j(p,T) \ge 0\) and that
\[
  \sigma_j(p, \lambda T) = |\lambda| \sigma_j(p,T)
\]
for all scalars \(\lambda \in \mathbb{R}\) (or \(\mathbb{C}\)).
Then for each fixed \(p\), the map
\[
  T \mapsto \|A_p T\|_{\mathrm{mult}}
\]
is positive and absolutely homogeneous:
\begin{enumerate}[leftmargin=1.5em]
  \item \(\|A_p T\|_{\mathrm{mult}} \ge 0\), with equality if and only if \(\sigma_j(p,T) = 0\) for all \(j\in J\),
  \item \(\|A_p(\lambda T)\|_{\mathrm{mult}} = |\lambda|\cdot \|A_p T\|_{\mathrm{mult}}\).
\end{enumerate}
\end{lemma}

\begin{proof}
By definition,
\[
  \|A_p T\|_{\mathrm{mult}} = \sum_{j\in J} \omega(p,\mu_j)\sigma_j(p,T)
\]
with \(\omega(p,\mu_j)\ge 0\) and \(\sigma_j(p,T)\ge 0\).
Thus \(\|A_p T\|_{\mathrm{mult}} \ge 0\).
If \(\|A_p T\|_{\mathrm{mult}} = 0\), then every term in the sum is non-negative and their sum is zero, implying \(\omega(p,\mu_j)\sigma_j(p,T) = 0\) for all \(j\).
If at least one \(\omega(p,\mu_j) > 0\), then \(\sigma_j(p,T) = 0\) for that \(j\); if all weights are zero, the functional is trivial and the lemma holds vacuously.
For homogeneity,
\[
  \|A_p(\lambda T)\|_{\mathrm{mult}}
  = \sum_{j\in J} \omega(p,\mu_j)\sigma_j(p,\lambda T)
  = \sum_{j\in J} \omega(p,\mu_j) |\lambda|\sigma_j(p,T)
  = |\lambda| \|A_p T\|_{\mathrm{mult}}.
\]
\end{proof}

\begin{remark}
No triangle inequality is assumed or required for the security arguments; \(\|\cdot\|_{\mathrm{mult}}\) plays the role of a structured observable rather than a norm in the strict functional-analytic sense.
\end{remark}

\section{Operator Norm Bounds in the Toy Model}

In the toy PITN+PWEH example, the tensor state is represented by vectors \(v(t) \in \mathbb{R}^3\) and the prime operators correspond to matrices
\[
  M_2 =
  \begin{pmatrix}
  1 & 1 & 0 \\
  0 & 1 & 1 \\
  0 & 0 & 1
  \end{pmatrix},\quad
  M_3 =
  \begin{pmatrix}
  1 & 0 & 1 \\
  1 & 1 & 0 \\
  0 & 1 & 1
  \end{pmatrix},\quad
  M_5 =
  \begin{pmatrix}
  2 & 0 & 0 \\
  0 & 1 & 1 \\
  1 & 0 & 1
  \end{pmatrix}.
\]

We work with the standard Euclidean norm \(\|\cdot\|_2\) on \(\mathbb{R}^3\).

\begin{lemma}[Bounds on \(\|M_p\|_2\)]
There exist constants \(C_2, C_3, C_5 > 0\) such that for all \(v \in \mathbb{R}^3\),
\[
  \|M_2 v\|_2 \le C_2 \|v\|_2,\quad
  \|M_3 v\|_2 \le C_3 \|v\|_2,\quad
  \|M_5 v\|_2 \le C_5 \|v\|_2.
\]
Moreover, one can take, for example,
\[
  C_2 = \sqrt{7},\quad C_3 = \sqrt{6},\quad C_5 = \sqrt{6}.
\]
\end{lemma}

\begin{proof}
We bound each operator via its Frobenius norm, using \(\|A\|_2 \le \|A\|_F\). [web:69]

Compute:
\[
  \|M_2\|_F^2 =
  1^2 + 1^2 + 0^2 + 0^2 + 1^2 + 1^2 + 0^2 + 0^2 + 1^2
  = 5,
\]
so \(\|M_2\|_2 \le \sqrt{5}\).
Similarly,
\[
  \|M_3\|_F^2 =
  1^2 + 0^2 + 1^2 + 1^2 + 1^2 + 0^2 + 0^2 + 1^2 + 1^2 = 6,
\]
\[
  \|M_5\|_F^2 =
  2^2 + 0^2 + 0^2 + 0^2 + 1^2 + 1^2 + 1^2 + 0^2 + 1^2 = 7.
\]
Thus
\[
  \|M_2\|_2 \le \sqrt{5},\quad
  \|M_3\|_2 \le \sqrt{6},\quad
  \|M_5\|_2 \le \sqrt{7}.
\]

Taking slightly looser but simpler constants, such as
\(C_2 = \sqrt{7}, C_3 = \sqrt{6}, C_5 = \sqrt{7}\), suffices for all \(v\), since \(\|M_p v\|_2 \le \|M_p\|_2\|v\|_2 \le C_p \|v\|_2\).
\end{proof}

\begin{corollary}[Growth bound along a prime sequence]
Let \(v(0) \in \mathbb{R}^3\) and a sequence of primes \((p_{i_t})_{t=1}^T\) be given.
Define \(v(t) = M_{p_{i_t}} v(t-1)\).
Then
\[
  \|v(T)\|_2 \le \Big(\prod_{t=1}^T C_{p_{i_t}}\Big) \|v(0)\|_2,
\]
where each \(C_{p_{i_t}}\) is a bound for \(\|M_{p_{i_t}}\|_2\) as above.
\end{corollary}

\begin{proof}
By induction using submultiplicativity:
\[
  \|v(T)\|_2 = \|M_{p_{i_T}} v(T-1)\|_2
  \le \|M_{p_{i_T}}\|_2 \|v(T-1)\|_2
  \le C_{p_{i_T}} \|v(T-1)\|_2,
\]
and iterating gives the product bound.
\end{proof}

\section{Non-Associativity and Order Sensitivity}

We now formalize the non-associative behavior in the toy model.

\begin{lemma}[Non-commutativity of operators]
In the toy PITN, the matrices \(M_2\) and \(M_3\) do not commute:
\[
  M_2 M_3 \neq M_3 M_2.
\]
\end{lemma}

\begin{proof}
Direct calculation:
\[
  M_2 M_3 =
  \begin{pmatrix}
  1 & 1 & 0 \\
  0 & 1 & 1 \\
  0 & 0 & 1
  \end{pmatrix}
  \begin{pmatrix}
  1 & 0 & 1 \\
  1 & 1 & 0 \\
  0 & 1 & 1
  \end{pmatrix}
  =
  \begin{pmatrix}
  2 & 1 & 1 \\
  1 & 2 & 1 \\
  0 & 1 & 1
  \end{pmatrix},
\]
\[
  M_3 M_2 =
  \begin{pmatrix}
  1 & 0 & 1 \\
  1 & 1 & 0 \\
  0 & 1 & 1
  \end{pmatrix}
  \begin{pmatrix}
  1 & 1 & 0 \\
  0 & 1 & 1 \\
  0 & 0 & 1
  \end{pmatrix}
  =
  \begin{pmatrix}
  1 & 1 & 1 \\
  1 & 2 & 1 \\
  0 & 1 & 2
  \end{pmatrix}.
\]
These matrices are not equal (for instance, the \((1,1)\) entry is \(2\) vs.\ \(1\)), hence \(M_2 M_3 \neq M_3 M_2\).
\end{proof}

\begin{corollary}[Order sensitivity of trajectories]
There exists \(v(0)\) such that
\[
  M_2 M_3 v(0) \neq M_3 M_2 v(0).
\]
\end{corollary}

\begin{proof}
Take \(v(0) = (1,1,1)^\top\).
Then
\[
  M_2 M_3 v(0) = (4,4,2)^\top,\quad
  M_3 M_2 v(0) = (3,4,3)^\top.
\]
Thus they differ.
\end{proof}

This order sensitivity is precisely what PWEH exploits: the sequence of primes induces a path through state space that cannot be freely reordered without changing observables.

\section{Path-Commitment and Collision Arguments}

\subsection{Modeling the Hash Chain}

We model the integrity chain as follows.
Let \(H : \{0,1\}^* \to \{0,1\}^n\) be a post-quantum hash function. [web:49][web:67]
For a sequence of steps indexed by \(t=1,\dots,T\), define:
\[
  S_0 = s_0 \in \{0,1\}^n,
\]
\[
  S_t = H\big(S_{t-1} \parallel X_t\big),
\]
where
\[
  X_t = \mathsf{enc}\big(p_{i_t}, \|A_{p_{i_t}}T(t)\|_{\mathrm{mult}}, M(t)\big)
\]
is a canonical encoding of the prime, the multiplicity-weighted norm, and metadata.

Thus, the pair \((S_t,X_t)\) encodes the hash chain and inputs at each step.

\subsection{Path-Collision Definition}

\begin{definition}[Path-collision]
Fix an integrity sequence \((S_t)_{t=0}^T\).
A \emph{path-collision} is a distinct sequence of inputs \((X'_t)_{t=1}^T\) such that:
\[
  S_0 = S'_0,\quad S'_t = H(S'_{t-1} \parallel X'_t)\ \ \text{for all }t,
\]
and
\[
  S'_t = S_t \quad\text{for all }t=0,\dots,T,
\]
but there exists at least one \(t\) with \(X'_t \neq X_t\).
\end{definition}

In the PWEH setting, any such path-collision must respect additional constraints:
\begin{itemize}[leftmargin=1.5em]
  \item Each \(X_t\) must arise from \((p_{i_t},T(t),M(t))\) under the PITN dynamics and policy rules,
  \item The same is true for \(X'_t\),
  \item Policy constraints forbid arbitrary primes and metadata.
\end{itemize}

\subsection{Informal Hardness Argument}

Under standard assumptions for \(H\), finding a path-collision in the unconstrained setting already implies the ability to find multi-step collisions in a hash chain.

More concretely, if an adversary can, for some fixed integrity sequence, construct an alternate input sequence \((X'_t)\) with all \(X'_t \neq X_t\) at some index but the same \((S_t)\), then they can be used to build non-trivial collisions for \(H\) itself. [web:68][web:72]

In PWEH, the adversary's job is harder:
each \(X_t\) encodes structured data (\(p_{i_t}\), \(\|A_{p_{i_t}}T(t)\|_{\mathrm{mult}}\), and \(M(t)\)) that must be consistent with non-associative PITN dynamics and policy manifolds.
Thus, the space of admissible input sequences is a strict subset of \(\{0,1\}^*\), and any path-collision must lie in this constrained space.

\begin{lemma}[Order-commitment in the toy model (informal)]
In the toy PITN+PWEH example, consider a sequence of primes \((2,3,2)\) producing integrity states \((S_0,S_1,S_2,S_3)\) from initial vector \(v(0)=(1,1,1)^\top\).
Assuming \(H\) is collision-resistant and \(\mathsf{enc}\) is injective on tuples, any distinct prime sequence of the same length that yields the same integrity sequence must violate the PITN dynamics or policy constraints.
\end{lemma}

\begin{proof}[Sketch]
For the sequence \((2,3,2)\), the toy model yields specific vectors:
\[
  v(1) = M_2 v(0), \quad v(2) = M_3 v(1), \quad v(3) = M_2 v(2),
\]
and associated norms \(\|A_{2}T(1)\|_{\mathrm{mult}},\ \|A_{3}T(2)\|_{\mathrm{mult}},\ \|A_2 T(3)\|_{\mathrm{mult}}\) as explicitly computed in the main text.

Suppose another prime sequence \((p'_{i_1},p'_{i_2},p'_{i_3})\) with corresponding states \((T'(1),T'(2),T'(3))\) and metadata \((M'(1),M'(2),M'(3))\) yields the same integrity chain \((S_1,S_2,S_3)\) while respecting policy and PITN rules.
Injectivity of \(\mathsf{enc}\) and collision-resistance of \(H\) then imply that at each step
\[
  X_t = X'_t,
\]
so
\[
  p_{i_t} = p'_{i_t},\quad \|A_{p_{i_t}}T(t)\|_{\mathrm{mult}} = \|A_{p'_{i_t}}T'(t)\|_{\mathrm{mult}},\quad M(t) = M'(t).
\]
In particular, the prime sequences coincide, which establishes order-commitment: no distinct prime order can yield the same chain without violating the hash assumptions or leaving the admissible PITN+policy space.
\end{proof}

\begin{remark}
A fully rigorous reduction to hash security would require formalizing the PITN state space and policy manifolds as an admissible input language for \(H\), then showing that an attacker that forges a path-collision there yields a collision-finder for \(H\).
The toy example illustrates the structure of such a reduction.
\end{remark}

\section{Policy-Enforced Bounds on Recursion}

Recall that policy manifolds \(\pi\) impose bounds such as
\[
  \beta(t) \le \beta_{\max}(t),\quad
  \mathcal{P}_{\mathrm{allowed}}(t) \subseteq \mathcal{P},
\]
and that these are encoded inside \(M_{\mathrm{policy}}(t)\).

\begin{definition}[Policy-compliant trace]
A trace \(\mathsf{Tr} = \{(p_{i_t}, T(t), M(t), S_t)\}_{t=0}^T\) is \emph{policy-compliant} under manifold \(\pi\) if for all \(t\):
\begin{enumerate}[leftmargin=1.5em]
  \item \(S_t = H(S_{t-1} \parallel X_t)\) with \(X_t\) derived from \((p_{i_t},T(t),M(t))\),
  \item \(M_{\mathrm{policy}}(t)\) equals \(\pi(t)\) or a prefix-consistent value,
  \item \(p_{i_t} \in \mathcal{P}_{\mathrm{allowed}}(t)\) as specified by \(\pi\),
  \item \(\beta(t) \le \beta_{\max}(t)\) as specified by \(\pi\).
\end{enumerate}
\end{definition}

\begin{lemma}[Policy-enforced recursion bound]
Let \(\pi\) be a policy manifold with recursion bound \(\beta_{\max}(t)\).
If an AGI trajectory \(\mathsf{Tr}\) achieves \(\beta(t) > \beta_{\max}(t)\) at any time \(t\), then \(\mathsf{Tr}\) is not policy-compliant under \(\pi\).
In particular, no verifier that requires compliance with \(\pi\) will accept the corresponding integrity state \(S_t\) as legitimate.
\end{lemma}

\begin{proof}
By definition of policy-compliance, condition (4) requires \(\beta(t) \le \beta_{\max}(t)\) at all times.
If for some time \(t^\ast\) we have \(\beta(t^\ast) > \beta_{\max}(t^\ast)\), then condition (4) is violated at \(t^\ast\), hence \(\mathsf{Tr}\) is not policy-compliant.
Any verifier insisting on \(\pi\) therefore rejects \(S_{t^\ast}\).
\end{proof}

\begin{corollary}[Bound on recognized intelligence scaling]
If access to external resources (compute, data, network) is conditioned on presenting a policy-compliant integrity state under \(\pi\), then no trajectory that exceeds \(\beta_{\max}\) is recognized as legitimate.
Thus the system cannot both exceed the recursion bound and maintain recognized, attested operation.
\end{corollary}

\section{Summary of Appendix Results}

In this appendix we have:
\begin{itemize}[leftmargin=1.5em]
  \item Formalized the multiplicity-weighted norm and established its basic properties,
  \item Derived explicit operator norm bounds for the toy PITN matrices \(M_2, M_3, M_5\),
  \item Demonstrated non-commutativity and order sensitivity for prime-indexed operators,
  \item Defined path-collision in the PWEH hash chain and argued informally that forging such a collision implies breaking the underlying post-quantum hash,
  \item Shown that policy manifolds impose hard, verifiable bounds on recursion depth and prime usage.
\end{itemize}

These results support the claim that PWEH provides a mathematically grounded, path-sensitive security and governance layer tightly coupled to the prime-structured dynamics of QAGI.

\nocite{*}
\bibliographystyle{unsrt}  
\bibliography{references}
\end{document}